% Drafted 2026-08-27 for the "Marijuana Spillover effect" Overleaf doc (main.tex).
% Revenue-per-local-resident dispensary-level check (1 km / 2 km denominator).
% NOT yet added to the paper. Uses \mathbbm{1} (bbm package, matches the doc).
% Figures to upload alongside:
%   fig_eventstudy_nm_dispensary_revper1kmresident.png
%   fig_eventstudy_nm_dispensary_revperresident_1km_vs_2km.png
% Source: analyze_eventstudy_nm_dispensary_revper1km.py, build_nm_dispensary_local_population.py

% ======================= BODY / APPENDIX PASSAGE =======================

\subsection{Revenue per local resident}
\label{sec:localresident}

The ZIP- and county-level per-capita results divide dispensary revenue by an
administrative population that may be a poor proxy for the customer base of an
individual store: two dispensaries in the same county, or even the same ZIP
code, can face very different resident populations within walking or short
driving distance. We therefore re-estimate the border premium at the
dispensary level using each store's \emph{own} local resident population as the
denominator.

For each of the 1{,}192 geocoded New Mexico dispensaries we sum 2020 Census
block population (TIGER/Line \texttt{tabblock20}, area-apportioned to the buffer)
within one kilometre of the dispensary's coordinates. This local population
averages roughly 3{,}100 residents (median 2{,}800) and is itself somewhat lower
around border dispensaries (2{,}340) than interior dispensaries (3{,}540),
consistent with border stores locating in smaller border towns. The border
classification is unchanged: a dispensary is border-exposed if its own location
lies within 75 miles of the Texas, Oklahoma, or Utah state line with at least
1{,}000 nonlegal-state residents within that reach.

The outcome is $\log(\text{own quarterly revenue} / \text{residents within 1 km})$.
We estimate a dynamic specification that interacts the border indicator with
every quarter,
%
\begin{equation}
\log\!\left(\frac{\text{Revenue}_{it}}{\text{Pop1km}_i}\right)
= \alpha + \sum_q \gamma_q\left(\text{Border}_i \times \mathbbm{1}[\text{Quarter}=q]\right)
+ \phi \log(\text{Pop1km}_i) + \delta_q + \eta_{b(i)} + \varepsilon_{it},
\end{equation}
%
where $\eta_{b(i)}$ is a chain/brand fixed effect and standard errors are
clustered by county. Each $\gamma_q$ is the border--interior difference in log
revenue per local resident in quarter $q$. Because a one-kilometre radius leaves
a thin tail of dispensaries with almost no nearby population, we drop the 37
stores with fewer than 100 residents within 1 km and winsorise the outcome at
its 1st and 99th percentiles; we also re-run the specification with a
two-kilometre denominator, which is less sensitive to geocoding error.

The border premium in revenue per local resident is positive but not
statistically distinguishable from zero. The pooled static estimate is
$+16.7$ percent ($p=0.286$; $n=9{,}918$ dispensary-quarters, 1{,}134
dispensaries), and $+14.9$ percent ($p=0.306$) with the two-kilometre
denominator. Figure~\ref{fig:es_revper1km} shows the dynamic estimates: the gap
is positive in every quarter from 2022Q3 onward, fluctuating between roughly
$+0.05$ and $+0.35$ log points with no trend, but only the 2024Q1 coefficient is
individually significant at the 5 percent level. Figure~\ref{fig:es_revper1km2km}
confirms that the one- and two-kilometre denominators trace essentially the same
path. This mirrors the other dispensary-level revenue results
(Section~\ref{sec:appendixdispensary}): once revenue is measured per store rather
than per market, the border advantage is small and imprecise. It is the number
of stores serving the border market, not the revenue of any one of them, that
responds to cross-border demand.

% ======================= FIGURE: MAIN (1 km) =======================

\begin{figure}[htbp]
\centering
\includegraphics[width=\textwidth]{fig_eventstudy_nm_dispensary_revper1kmresident.png}
\caption{Dynamic border--interior difference in dispensary revenue per local
resident, New Mexico. Each point is the coefficient on
$\text{Border}_i \times \mathbbm{1}[\text{Quarter}=q]$ from a dispensary-quarter
regression of $\log(\text{own revenue}/\text{residents within 1 km})$ on the full
set of border-by-quarter interactions, a log local-population control, quarter
fixed effects, and chain/brand fixed effects; county-clustered 95 percent
confidence intervals. Local resident population is 2020 Census block population
(area-apportioned) within one kilometre of each dispensary's geocoded location.
Sample: dispensary-quarters with positive revenue and at least 100 residents
within 1 km, outcome winsorised at the 1st/99th percentiles (9{,}918
dispensary-quarters, 1{,}134 dispensaries, 2022Q2--2026Q2; 2026Q3 dropped as a
partial quarter). The pooled static border coefficient is $+16.7$ percent
($p=0.286$).}
\label{fig:es_revper1km}
\end{figure}

% ======================= FIGURE: ROBUSTNESS (1 km vs 2 km) =======================

\begin{figure}[htbp]
\centering
\includegraphics[width=\textwidth]{fig_eventstudy_nm_dispensary_revperresident_1km_vs_2km.png}
\caption{Revenue per local resident, one- versus two-kilometre denominator.
Border-by-quarter coefficients from the specification of
Figure~\ref{fig:es_revper1km}, with the resident-population denominator computed
over a one-kilometre radius (as in the main specification) and a two-kilometre
radius (less sensitive to geocoding error; the two-kilometre sample drops
dispensaries with fewer than 250 residents in that radius). County-clustered
95 percent confidence intervals. The two denominators yield nearly identical
dynamic paths; the pooled static border coefficient is $+16.7$ percent
($p=0.286$) at one kilometre and $+14.9$ percent ($p=0.306$) at two kilometres.}
\label{fig:es_revper1km2km}
\end{figure}

% ======================= 100-WORD INTERPRETATION (for Results/Discussion) =======================

% Measuring each dispensary's revenue against the actual resident population within
% one kilometre -- its true local customer base, not an administrative aggregate --
% the border premium in per-resident revenue is small and statistically
% indistinguishable from zero (+16.7%, p=0.286; +14.9%, p=0.306 at two kilometres),
% with a flat, trendless dynamic path. This is the most demanding test of the
% per-firm channel, and it fails. That failure strengthens the primary results: the
% large border premium in revenue per capita and dispensary density is not a
% denominator artifact, and it does not operate by enriching incumbent stores.
% Cross-border demand is absorbed through entry -- more dispensaries -- leaving
% per-store revenue near parity.
